\subsection{Step 3: Construction of the Manual Review Table}
\label{sec:step3_review_table}

Step~3 converts the Step~2 candidate-screening output into a manual-review table. This step is intentionally a review stage, not a merge stage. It does not assign canonical names, create harmonized groups, or collapse the node universe. Its sole purpose is to identify which candidate pairs should be examined before any equivalence decision is propagated downstream.

The direct inputs are \texttt{02\_candidate\_pairwise\_similarity.csv}, \texttt{02\_screening\_cutoff\_selection.csv}, and \texttt{02\_exact\_name\_candidate\_stratum.csv}. Candidate pairs are retained if they satisfy at least one of the provisional screening cutoffs. Exact-name pairs, if any exist, remain a separate review stratum rather than an automatic merge rule. The review template also carries forward whether either textual instance was flagged in Step~1 as requiring manual provenance validation.

The corrected design restores the three allowed review states:
\begin{itemize}
    \item \texttt{equivalent}
    \item \texttt{not\_equivalent}
    \item \texttt{uncertain}
\end{itemize}
The columns \texttt{review\_decision} and \texttt{review\_notes} are left blank by default. No automatic equivalence-by-cutoff rule is applied.

\subsubsection{Outputs of this Step}
\label{sec:step3_outputs}

The corrected Step~3 produces:
\begin{itemize}
    \item \texttt{03\_review\_pairs\_textual\_instances.csv}
    \item \texttt{03\_review\_summary.csv}
    \item \texttt{03\_review\_template\_textual\_instances.csv}
\end{itemize}

The review template is the file to be completed by the analyst before true harmonization can begin. It retains the similarity evidence from Step~2 together with the empty review fields and any provenance warning indicators inherited from Step~1.

\subsubsection{Fields to Fill After Concluding This Step}
\label{sec:step3_fields_to_fill}

The following columns in \texttt{03\_review\_template\_textual\_instances.csv} are meant to be completed after inspection:
\begin{itemize}
    \item \texttt{review\_decision}
    \item \texttt{review\_notes}
\end{itemize}

Until those fields are filled, Step~4 should remain in a no-merge baseline.

\subsubsection{Implementation Note and Current Run}
\label{sec:step3_implementation_note}

The current corrected run retained:
\begin{itemize}
    \item 160 total review pairs;
    \item 144 pairs passing the \texttt{name\_only} screening cutoff;
    \item 57 pairs passing the \texttt{name\_description} screening cutoff;
    \item 41 pairs passing both screening cutoffs;
    \item 16 review pairs involving at least one text instance flagged for manual provenance validation;
    \item 0 exact-name review-stratum pairs.
\end{itemize}

Most importantly, the current review file is unfilled:
\begin{itemize}
    \item \texttt{review\_decisions\_pre\_filled = false}
    \item \texttt{allowed\_review\_states = equivalent;not\_equivalent;uncertain}
\end{itemize}

This means Step~3 is now aligned with the intended methodology. The workflow pauses here for actual judgment rather than converting screening cutoffs directly into harmonization decisions.
